\section{Vertex cover}
\label{sec:vc}

In this section we consider the parameterization by the size of a \emph{vertex cover}, which is a set $S\subseteq V$ of vertices such that every edge is incident on at least one of the vertices of $S$. 
We first present an easy FPT algorithm based on the dynamic program given by 
\cref{thm:tw-dp}, and then prove that its runtime dependence on the parameter is 
asymptotically optimal.
%

%
\subsection{FPT Algorithm}

Let $S\subseteq V$ be a given vertex cover of size $k$ for the input graph $G$. We will show how to construct a tree decomposition with all required properties of \cref{thm:tw-dp} in order to run the corresponding dynamic program. For this the tree decomposition $(T,\{B_i\}_{i\in V(T)})$ needs to be \emph{nice}.

We may assume w.l.o.g.\ that the vertex cover $S$ contains no terminal: using a standard preprocessing procedure, we can replace any terminal $t\in S$ of the vertex cover by a Steiner vertex $v$ and then connect $t$ with $v$ using an edge of cost~$0$. 
Note that $S$ is still a vertex cover for the preprocessed graph and that the complement set $I=V\setminus S$
of the vertex cover is an independent set containing all terminals. To use \cref{thm:tw-dp},
we will first construct a (trivial) nice tree decomposition for~$I$ and then add $S$ to each bag.

Note that a terminal $t\in R$ can be part of several demand pairs of the \SF instance. 
Consider the \emph{demand graph} $H$ with vertex set $R$ and an edge for each demand pair.
Any subset of $R$ that induces a maximal connected component of $H$ is called a \emph{group}.
Note that every group of terminals must lie in the same connected component of any \SF solution.
For each group $R'\subseteq R$, we create one leaf node of the tree decomposition for each terminal $t\in R'$
and let the corresponding bag contain $t$. We then add a forget node for each such leaf node, which we add
as parent to the leaf with an empty bag. These forget nodes are then connected in a binary tree by adding
join nodes with empty bags (unless the group only contains one terminal in which case we skip this step). 
We proceed in the same way for the Steiner vertices of the independent set $I$,
that is, if we consider $I\setminus R$ to be a group as well we obtain a nice tree decomposition for $I\setminus R$
in which each bag of a leaf node contains one vertex of $I\setminus R$. 
All these trees are then connected using join nodes with empty bags, to obtain a tree decomposition
(of width $0$) for the independent set $I$. Finally, we simply add the vertex cover $S$ to every bag,
which results in a nice tree decomposition (of width $k$) for the graph $G$, such that
every terminal lies in a bag of a leaf node (as~$S\cap R=\emptyset$).

In the obtained tree decomposition, let $V_B$ be the vertices of $G$ contained in all bags in the subtree rooted at the node associated with $B$. By construction, $V_B$ either contains no terminals 
(if $B$ is a bag of the tree decomposition for $I\setminus R$), 
fully contains some groups of $R$ (if $B$ is the bag of the root of a tree decomposition for a group $R'\subseteq R$, 
or if $B$ is a bag of a join node used to connect the tree decompositions for groups in the last step),
or contains some strict subset of only one group of~$R$
(if $B$ is a bag of a non-root node of a tree decomposition for a group~$R'\subseteq R$). 
If no terminals lie in $V_B$ then clearly there are no active terminals for bag $B$. 
However this is also the case if $V_B$ fully contains some groups of~$R$. Hence in both these cases
the set $\Pi_B$ of permutations of active terminals is empty. 
Whenever~$V_B$ contains a strict subset of only one group $R'\subseteq R$,
the active terminals $A_B$ of $B$ are only from this set, i.e., $A_B\subseteq R'$. Thus we can
add the trivial partition $\pi=\{A_B\}$ as the only partition of $\Pi_B$, since all terminals of $R'$ belong to the same component of any solution, including the optimum. 

Clearly, the optimal solution conforms with these sets $\Pi_B$ of permutations, and the total
number $p$ of permutations is at most the number of groups, which is at most $n/2$.
Hence by \autoref{thm:tw-dp} we obtain the algorithm of \autoref{thm:vc}.

%
